\section*{Chapitre \ref{ch:in:revenue-and-profit-across-the-industry} --- Revenus et bénéfices dans l'ensemble du secteur}

\begin{solution}{exo:in:revenue-and-profit-across-the-industry:1}
Les revenus moins les charges liées aux transactions sont égaux aux charges d'exploitation plus le résultat d'exploitation, soit $5\,002+4\,929=\$9\,931$ millions~;
les charges liées aux transactions s'élèvent à $12\,640-9\,931=\$2\,709$ millions.
\end{solution}

\begin{solution}{exo:in:revenue-and-profit-across-the-industry:2}
Nasdaq (0.55~million de dollars), Man Group et l'Intercontinental Exchange (0.77~million de dollars chacun), Morgan Stanley (0.85
million de dollars), Flow Traders (0.86~million de dollars) et MarketAxess (0.94~million de dollars). Le chiffre le plus élevé pour une société entière, celui de Quadrature
9.29~millions de dollars, soit 16.9 fois celui de Nasdaq.
\end{solution}

\begin{solution}{exo:in:revenue-and-profit-across-the-industry:3}
CME Group~: $4\,229.5/3\,875=\$1.09$~million. Optiver~: $1\,369/2\,100=\text{\euro}0.65$~million.
\end{solution}

\begin{solution}{exo:in:revenue-and-profit-across-the-industry:4}
Son effectif est passé de 6\,377 à 8\,525 au cours de 2023, «~principalement en raison de notre acquisition d'Adenza~»~: l'effectif acquis
est compté en totalité à la fin de l'année, alors que le revenu acquis n'arrive que sur les années suivantes.
\end{solution}

\begin{solution}{exo:in:revenue-and-profit-across-the-industry:5}
$(29.25/16.85)^{1.19}=1.92$~: le \omterm{def:in:revenue-and-profit-across-the-industry:perhead}{revenu par tête} est environ 92~\% plus élevé. Le rapport réel de Virtu était de 2.16 et celui de Flow Traders de 4.51~:
l'élasticité ajustée est la moyenne de deux sociétés dont la sensibilité diffère, et la moyenne annuelle du VIX n'est qu'un parmi plusieurs éléments qui
fixent le revenu d'un teneur de marché.
\end{solution}

\begin{solution}{exo:in:revenue-and-profit-across-the-industry:6}
C'est une filiale dont le revenu inclut des revenus de service fixés par le prix de transfert du groupe, et dont
l'effectif est passé de 96 à 21 lorsque le travail a été transféré ailleurs dans le groupe~: son \omterm{def:in:revenue-and-profit-across-the-industry:perhead}{revenu par tête} (de 1.00~million de dollars à 7.18~million de dollars) décrit les
comptes de l'entité, et non les salariés de la société.
\end{solution}

\begin{solution}{exo:in:revenue-and-profit-across-the-industry:7}
Les élasticités des teneurs de marché restent à 1.19 et celle des bourses à 0.13~; leurs erreurs types ne changent qu'à la troisième
décimale. Avec les effets fixes par société et une pente par catégorie, la pente de chaque catégorie est estimée à partir de
la variation propre à ses sociétés~; les autres catégories n'interviennent que par la variance résiduelle regroupée.
\end{solution}

\begin{solution}{exo:in:revenue-and-profit-across-the-industry:8}
Le haut du tableau est celui des sociétés qui déposent des comptes, et non celui du secteur~: les plus grands teneurs de marché privés ne publient aucun
résultat consolidé, de sorte que le \omterm{def:in:revenue-and-profit-across-the-industry:disclosure}{biais de publication} les écarte. Le montant de 13~millions de dollars est aussi une autre ligne (les commissions d'un gestionnaire de fonds)
que le revenu net de trading d'un teneur de marché, et elle provient de 98 personnes au cours d'une année exceptionnelle.
\end{solution}

\begin{solution}{pb:in:revenue-and-profit-across-the-industry:1}
\begin{enumerate}
\item Le revenu (après les coûts de pass-through) ou une ligne de résultat publiée, divisé par l'effectif~; avec la ligne, la base de comptage de l'effectif,
la devise et le niveau de prix.
\item Une ligne brute compte les sommes collectées pour le compte d'autrui~; le total des revenus du groupe boursier, 12\,640 millions de dollars, comprend
2\,709 millions de dollars de charges liées aux transactions, ce qui laisse 9\,931 millions de dollars.
\item Fin d'année, moyenne, équivalents temps plein, milliers arrondis, bornes inférieures indiquées~; les effectifs de fin d'année surestiment le
dénominateur d'une société en croissance, et une borne inférieure donne une borne supérieure par tête.
\item Taux de change annuels moyens de la BCE vers le dollar, puis indice américain des prix à la consommation pour exprimer le tout en dollars de 2025.
\item Quinze sociétés entières (plus une entité britannique présentée à part), 87 années-société, à partir des company facts de la SEC, des tableaux du livre~16,
des rapports annuels, des communiqués de résultats et du registre britannique.
\item Le plus bas~: Nasdaq, 0.48~million de dollars en 2023. Le plus haut~: Quadrature, 13.0~millions de dollars pour son exercice clos en janvier 2022.
\item Flow Traders 4.5, Virtu 2.5, CME Group 1.2.
\item Acquisitions~: l'effectif de Nasdaq est passé de 6\,377 à 8\,525 avec Adenza~; celui de l'Intercontinental Exchange est passé de
8\,911 à 13\,222.
\item Les banques déclarent un revenu de marchés mais aucun effectif par division.
\item Il est passé de 1.00~million de dollars à 7.18~million de dollars tandis que son \omterm{def:in:reading-a-trading-firms-accounts:staff}{effectif moyen} tombait de 96 à 21~; son revenu est fixé au sein
du groupe.
\item 15.39, 29.25, 19.66, 25.64, 16.85, 15.55, 18.93.
\item $\log(R/N)_{it}=a_i+b_{k(i)}\log\bar V_t+\varepsilon_{it}$~; les effets fixes par société retirent le niveau propre à chaque société, de sorte que les
pentes proviennent des variations à l'intérieur des sociétés.
\item Teneurs de marché 1.19 (0.26)~; bourses et plateformes 0.13 (0.16)~; sociétés de gestion 0.23 (0.27)~; banques 0.03 (0.32)~;
courtier $-0.39$ (0.37)~; bourse crypto $-1.27$ (0.59).
\item Son revenu suit les cours des cryptos, qui ont le plus baissé en 2022, une année de forte volatilité des actions.
\item Virtu 2.16, Flow Traders 4.51.
\item Teneurs de marché 1.19 (0.26), bourses et plateformes 0.13 (0.16), sociétés de gestion 0.23 (0.27), banques 0.03 (0.32),
courtier $-0.39$ (0.37), bourse crypto $-1.27$ (0.59)~; le rapport 2020/2023 du \omterm{def:in:revenue-and-profit-across-the-industry:perhead}{revenu par tête} va de 2.16 à 4.51
pour les teneurs de marché.
\item Les résidus partagent des années et ne sont pas indépendants~; sept années comportent un grand épisode de volatilité.
\item Le biais vient de ce que les sociétés qui publient diffèrent de celles qui ne publient pas~; il laisse de côté les sociétés privées les plus
rentables, de sorte qu'il sous-estime le haut du \omterm{def:in:bank-markets-divisions:league}{classement}.
\item Vers le haut~: les sociétés défaillantes sortent avec leurs mauvaises années.
\item Oui~: le teneur de marché a de bonnes années~; son revenu par personne a à peu près doublé ou davantage de 2023 à 2020, et il suit le marché bien plus qu'une
bourse. Celui de la bourse est plus stable, et la rémunération variable du teneur de marché le sera aussi
(chapitre~13).
\end{enumerate}
\end{solution}

\begin{solution}{iq:in:revenue-and-profit-across-the-industry:1}
Davantage de volume (plus de transactions sur lesquelles gagner une fourchette) et des fourchettes plus larges (davantage gagné par transaction, à mesure que le risque d'inventaire augmente)~; aussi
plus de dislocations entre instruments liés.
\iqlookfor{volume et fourchette, avec la raison du risque d'inventaire pour la fourchette.}
\end{solution}

\begin{solution}{iq:in:revenue-and-profit-across-the-industry:2}
Les revenus diminués des charges liées aux transactions, car le reste est collecté pour le compte d'autres (frais réglementaires) ou reversé aux
fournisseurs de liquidité~; les revenus nets d'une banque sont eux aussi calculés après les charges d'intérêts.
\iqlookfor{retrancher les coûts de pass-through avant de comparer.}
\end{solution}

\begin{solution}{iq:in:revenue-and-profit-across-the-industry:3}
L'effet fixe par société retire le niveau permanent de chaque société, de sorte que la pente provient des variations à l'intérieur des sociétés. Un effet annuel absorberait
tout choc annuel commun, y compris le VIX lui-même, qui ne varie que d'une année à l'autre~: la pente ne serait plus
identifiée, sauf comme différences entre catégories.
\iqlookfor{identification~: un régresseur au niveau annuel ne peut coexister avec des effets annuels.}
\end{solution}

\begin{solution}{iq:in:revenue-and-profit-across-the-industry:4}
Grouper les erreurs par année (le choc est commun), mais sept groupes sont trop peu nombreux pour la formule usuelle~; utiliser un wild bootstrap par
année, ou présenter l'intervalle des pentes société par société.
\iqlookfor{peu de groupes, et une méthode qui en tient compte.}
\end{solution}

\begin{solution}{iq:in:revenue-and-profit-across-the-industry:5}
Indexer chaque donnée par (déclarant, concept, début de période, fin de période, dépôt)~; conserver la valeur de chaque dépôt afin que les retraitements soient de nouvelles
lignes~; une vue retient le dernier dépôt par période. Stocker l'unité et la devise avec la valeur.
\iqlookfor{clés bitemporelles~: période et dépôt.}
\end{solution}

\begin{solution}{iq:in:revenue-and-profit-across-the-industry:6}
À partir d'éléments publics seulement~: l'effectif tiré de leurs propres comptes (en fourchettes), les fourchettes de revenu issues du \omterm{def:in:revenue-and-profit-across-the-industry:perhead}{revenu par tête} des pairs cotés
ajustées par toute mesure publique d'activité, et présentées comme des fourchettes assorties de leurs hypothèses. On refuserait
de les présenter comme un chiffre, de classer les sociétés selon elles, ou d'utiliser des chiffres de presse comme données.
\iqlookfor{fourchettes, hypothèses énoncées, et refus de la fausse précision.}
\end{solution}
